% Viscosity lab report. Build from ../ with `make` (runs the analysis first).
% Rules: ../../report-rules.md (general), ../report-rules.md (this experiment).
% Limits (Moodle section "צמיגות", 2026-11-21): 10 pages, 18 numbered formulas. Cover + appendices
% excluded; theory formulas (1)-(6) excluded from the 18 (G p.1, conflict C2).
% Preamble = the free-fall and heat-conduction reports' (LuaLaTeX + babel, Times New Roman).
\documentclass[10pt,a4paper]{article}

\usepackage[a4paper,margin=2cm]{geometry}
\usepackage{amsmath}
\usepackage{graphicx}
\usepackage{float}  % [H]: figures and tables stay where they are written (user, 2026-10-30)
\usepackage{tikz}
\usepackage{xcolor}
\usepackage[labelsep=colon,font=small]{caption}
\usepackage{fontspec}
% LuaLaTeX + babel instead of XeLaTeX + polyglossia (2026-10-28): under XeTeX's bidi,
% \textcolor leaked to whole RTL lines; luacolor colours per character.
\usepackage{luacolor}
\usepackage[bidi=basic, layout=lists.tabular, main=hebrew, english]{babel}
% Clickable references to the report's own formulas/figures. hidelinks: no box, no colour change, so a
% reference keeps the font and colour of the text around it (user, 2026-11-22). Booklet formulas ("בחוברת") stay plain text.
\usepackage[hidelinks]{hyperref}
% Times New Roman (macOS system font, has Hebrew). Frank Ruehl CLM's digits read as bold.
\babelfont[hebrew]{rm}{Times New Roman}
\newcommand{\textenglish}[1]{\foreignlanguage{english}{#1}}
% Second-level list labels (א), (ב)... In running text a sub-item is referred to as 1א (user, 2026-10-30).
\renewcommand{\theenumii}{\localecounter{letters}{enumii}}
\renewcommand{\labelenumii}{(\theenumii)}

% Paragraphs: no first-line indent, small gap between paragraphs (Word-like, as in the guide).
\setlength{\parindent}{0pt}
\setlength{\parskip}{4pt plus 1pt}

% Review colours (general rules §2.4c-d): blue = text Claude wrote, until approved; green = edits on
% the partner's suggestions, until approved. Off for submission: \reviewfalse turns both black,
% the macros stay so the source still records who wrote what.
\newif\ifreview \reviewtrue
\ifreview
  \colorlet{addedblue}{blue}
  \definecolor{partnergreen}{RGB}{0,140,60}
\else
  \colorlet{addedblue}{black}
  \colorlet{partnergreen}{black}
\fi
\newcommand{\added}[1]{\textcolor{addedblue}{#1}}
\newcommand{\partner}[1]{\textcolor{partnergreen}{#1}}

% Visible placeholder for anything not written yet. Must not survive into the submitted PDF.
\newcommand{\TODO}[1]{\textcolor{red}{[#1]}}
% A spot in verbatim staff text waiting for the user's decision (red, no brackets). Must not survive either.
\newcommand{\flagged}[1]{\textcolor{red}{#1}}
% Units after a value in math mode: $5.30\U{W}$
\newcommand{\U}[1]{\,\mathrm{#1}}

\InputIfFileExists{details.tex}{}{% Cover-page placeholders, tracked. Real values go in details.tex (gitignored: personal IDs, public repo).
\newcommand{\ExperimentName}{צמיגות}
\newcommand{\ExperimentPart}{\TODO{חלק}}
\newcommand{\InstructorName}{\TODO{שם מלא של המדריך}}
\newcommand{\ExperimentDate}{\TODO{תאריך ביצוע}}
\newcommand{\SubmissionDate}{\TODO{תאריך הגשה}}
\newcommand{\StudentOneFirst}{\TODO{שם פרטי}}
\newcommand{\StudentOneLast}{\TODO{שם משפחה}}
\newcommand{\StudentOneID}{\TODO{ת.ז.}}
\newcommand{\StudentTwoFirst}{\TODO{שם פרטי}}
\newcommand{\StudentTwoLast}{\TODO{שם משפחה}}
\newcommand{\StudentTwoID}{\TODO{ת.ז.}}
\newcommand{\StudyTrack}{\TODO{מסלול}}
\newcommand{\LabGroup}{\TODO{קבוצה}}
\newcommand{\SubGroup}{\TODO{תת קבוצה}}
\newcommand{\Station}{\TODO{עמדה}}
}
% Every reported number comes from analysis/analysis.py (make).
% Generated by analysis/analysis.py -- do not edit by hand.
% Values are in math mode; units are written in the report text.


\begin{document}
% Cover page, laid out after ../../cover-page-template.docx (Moodle template).
% Not counted in the page limit and not numbered.
\begin{titlepage}
\thispagestyle{empty}
\noindent
\fbox{\begin{minipage}[t][3.2cm]{0.52\textwidth}
  \vspace{0.4cm}
  \Large\textbf{דו"ח מסכם בניסוי: \ExperimentName}\\[0.6cm]
  \textbf{חלק: \ExperimentPart}
\end{minipage}}
\hfill
\fbox{\begin{minipage}[t][3.2cm]{0.36\textwidth}
  \vspace{0.3cm}
  שם הבודק: \rule{2.8cm}{0.4pt}\\[0.35cm]
  תאריך הבדיקה: \rule{2.4cm}{0.4pt}\\[0.35cm]
  ציון הדו"ח: \quad I \rule{1cm}{0.4pt}\\[0.2cm]
  \hspace*{1.95cm} II \rule{1cm}{0.4pt}
\end{minipage}}

\vspace{1.5cm}
\noindent שם מדריך הניסוי (שם מלא): \InstructorName\\[0.5cm]
תאריך ביצוע הניסוי: \ExperimentDate\\[0.5cm]
תאריך הגשת הדו"ח: \SubmissionDate

\vspace{1.2cm}
\noindent\textbf{הדו"ח מוגש על ידי:}

\vspace{0.6cm}
\noindent
\begin{tabular}{c@{\hspace{0.6cm}}c@{\hspace{0.6cm}}c@{\hspace{0.6cm}}c}
 & שם פרטי & משפחה & ת.ז. \\ \hline
I & \StudentOneFirst & \StudentOneLast & \StudentOneID \\[0.3cm]
II & \StudentTwoFirst & \StudentTwoLast & \StudentTwoID \\
\end{tabular}

\vspace{1cm}
\noindent
\begin{tabular}{cccc}
\StudyTrack & \LabGroup & \SubGroup & \Station \\ \hline
מסלול הלימוד & מס' קבוצת המעבדה & תת קבוצה & מספר עמדה \\
\end{tabular}

\vspace{1.5cm}
\noindent\rule{\textwidth}{1.2pt}

\vspace{0.4cm}
\noindent\underline{\textbf{הערות הבודק לנושאים לקויים בדו"ח:}}
\end{titlepage}

\setcounter{page}{1}

\section*{מטרות הניסוי}
% Staff goal (briefing p.1): "מציאת ערך צמיגות הגליצרין ואימות חוק סטוקס". Turned into falsifiable goals
% (general rules §4.1): one sentence per goal. Claude's wording, approved by the user 2026-11-22.
בניסוי זה נמצא את ערך צמיגות הגליצרין ונאמת את חוק סטוקס:
\begin{enumerate}
  \item נמצא את צמיגות הגליצרין בטמפרטורת החדר משתי שיטות התאמה, ונראה שכל אחד מהערכים תואם לערך התיאורטי בטווח הטעות המקובל.
  \item נאמת את חוק סטוקס בכך שנראה שמהירות הנפילה הסופית של כדור בגליצרין תלויה לינארית בריבוע רדיוס הכדור, כפי שהחוק צופה.
\end{enumerate}

\section*{רקע תיאורטי}
% Verbatim from ../theory-background-format.docx (F, the staff's theory-format file), 2026-11-22.
% Text unchanged except whitespace. Math retyped from the docx's equations, same numbering (1)-(6).
% F's formulas are correct; the briefing's typos (m and the buoyancy term, missing /3) are not in F
% (Ran: use what should be, 2026-11-21).
% Typos in F kept verbatim (user, 2026-11-22: "keep theory as is").
נוזל זורם המקיים את חוק ברנולי נקרא נוזל אידיאלי. במצב זה מהירויות כל חלקיקי הנוזל בחתך מסוים של צינור שוות; כלומר, הנוזל מתקדם כגוף אחד לאורך הצינור. בזרימת נוזל ריאלי מופיעים כוחות נוספים שלא מאפשרים את תיאור התנועה בעזרת חוק ברנולי. כוחות אלה הם כוחות חיכוך פנימיים המתנגדים לתנועת הנוזל, הם אופייניים לסוג הנוזל ומהווים את הצמיגות. במערכת היחידות \textenglish{c.g.s} נמדדת הצמיגות $\eta$ ביחידות הנקראת \textenglish{poise} כאשר:
\[
\mathrm{poise} = \frac{\mathrm{gr}}{\mathrm{cm} \cdot \mathrm{sec}}
\]
הכוחות הפועלים על כדור בתוך זורם: הכוח המעכב $R$, הפועל על כדור הנע בתוך זורם צמיג, לא דחיס ולמינארי, נתון לפי חוק סטוקס על ידי הנוסחה:
\begin{equation}\label{eq:stokes}
R = 6\pi\eta r u
\end{equation}
כאשר $r$ רדיוס הכדור, $u$ מהירות הכדור ו-$\eta$ מקדם הצמיגות של הנוזל. נזכור כי חוק סטוקס נכון רק לגבי זורם בעל ממדים אינסופיים ביחס לממדי הגוף הנע, ומהירויות נמוכות מהמהירות הקריטיות. במהירויות העולות על המהירות הקריטית נגרמת טורבולנציה, נוצרות מערבולות בזמן התנועה והחוק מאבד את תוקפו. חוק סטוקס מאפשר מדידה של מקדם צמיגות של זורם. את משוואת התנועה של כדור נופל תחת כוח הכובד בתוך זורם נקבל לפי החוק השני של ניוטון:
\begin{equation}\label{eq:newton}
ma = mg - R - B
\end{equation}
כאשר $B$ הוא כוח העילוי לפי ארכימדס, $g$ תאוצת הכובד ו-$m$ מסת הכדור.

כוח סטוקס מתחיל לפעול מיד עם תחילת תנועת הכדור בתוך הזורם, והוא הולך וגדל ביחס ישר למהירות. ברגע בו התנועה הופכת לתנועה במהירות קבועה מקבלים $0 = mg - R - B$ ונשארים עם איבר מהירות שנסמן ב-$v_f$ שייצג את מהירות הנפילה הסופית:
\begin{equation}\label{eq:vf}
v_f = \frac{2g(\rho_b - \rho_l)r^2}{9\eta}
\end{equation}
כאשר $\rho_b$ צפיפות הכדור, $\rho_l$ צפיפות הנוזל, $r$ רדיוס הכדור ו-$v_f$ מהירות סופית.

משוואה \eqref{eq:newton} היא למעשה משוואה דיפרנציאלית עבור מהירות הכדור. מפתרון המשוואה מתקבלת מהירות הכדור כתלות בזמן:
\begin{equation}\label{eq:vt}
v(t) = v_f + (v_0 - v_f)\exp\left(-\frac{t}{\tau}\right), \qquad \tau = \frac{2\rho_b r^2}{9\eta}
\end{equation}
כאשר $v_0$ היא מהירותו ההתחלתית של הכדור.

את צמיגותו של נוזל ניתן להגדיר כהתנגדות שלו לשינוי צורתו שמקורה בחיכוך שבין חלקיקי הנוזל. העלאת הטמפרטורה של הנוזל גורמת מצד אחד לעליה באנרגיה הקינטית של חלקיקי הנוזל, ולכן לעליה במספר ההתנגשויות ביניהם ולעליה בחיכוך, ומצד שני להקטנת צפיפות הנוזל, ולכן להקטנת החיכוך. באופן ניסיוני נמצא שהגורם השני משמעותי יותר ולכן הצמיגות יורדת ככל שהטמפרטורה עולה.
\begin{equation}\label{eq:etaT}
\eta = A e^{\frac{B}{T}}
\end{equation}
כאשר $\eta$ הינה הצמיגות ביחידות של \textenglish{poise}, $T$ הטמפרטורה ביחידות צלזיוס ו-$A, B$ הינם קבועים אמפיריים הנקבעים לפי גרף נסוני של צמיגיות כתלות בטמפרטורה. מנוסחה \eqref{eq:etaT} ומיחידות הצמיגות אנו יכולים להסיק כי הצמיגות תלויה בזמן הנפילה ולכן נקבל את הקשר:
\begin{equation}\label{eq:tT}
t = C e^{\frac{D}{T}}
\end{equation}
כאשר $t$ הינו זמן נפילת כדור ביחידות של שניות, $T$ הטמפרטורה ביחידות צלזיוס ו-$C, D$ הינם קבועים אמיפריים הנקבעים לפי גרף ניסיוני של זמן כתלות בטמפרטורה בעלי יחידות שניות ומעלות צלזיוס.

\section*{מהלך הניסוי}
% Two components (guide p. 2-3): the system (description, equipment list with every instrument's
% resolution, schematic figure numbered like the list) and the procedure (reproducible from this
% section alone). Sources: briefing B (equipment p.1, stages p.3-6, appendix ב), the data file
% (resolutions, frame times, row counts), the user's answers 2026-11-21 (viscosity/todos.md 4a-4f).
% Everything here is Claude's wording: blue until the user approves (written 2026-11-22).
\subsection*{מערכת הניסוי}
מערכת הניסוי מתוארת באיור~\ref{fig:system}. מרכזה הוא מבחנה אנכית מלאה בגליצרין, חתומה בשסתום, שלאורכה סמני גובה שניתן להזיז. כדורים בשישה קטרים שונים מופלים לתוך הגליצרין מפתח המבחנה. מול המבחנה מוצבת על חצובה מצלמת סמארטפון המצלמת את נפילת הכדור, ומאחורי המבחנה בריסטול לבן המשמש כרקע. הסרטונים מנותחים בתוכנת העקיבה \textenglish{Tracker}, הנותנת את מיקום הכדור בכל פריים. צפיפות הגליצרין נמדדת באריומטר, וטמפרטורת הגליצרין נמדדת בחיישן טמפרטורה המחובר למחשב עם תוכנת \textenglish{Capstone}.

\subsection*{רשימת ציוד}
% Numbers match the labels in figure 1. Resolutions: data file (part a sheet); ruler 1 mm (user, 2026-11-22);
% temperature 0.001 C = last digit of the reading 23.116 (user: infer from the reading, 2026-11-22). Frame interval:
% Tracker time column, 0.0333 s steps = 30 fps (briefing p.1).
\begin{enumerate}\setlength{\itemsep}{0pt}\setlength{\parskip}{0pt}
  \item מבחנה מלאה בגליצרין, חתומה בשסתום.
  \item סמני גובה ניתנים להזזה על המבחנה.
  \item כדורים בשישה קטרים שונים, כ-6 עד 10 כדורים מכל קוטר, במגש.
  \item מאזניים דיגיטליים, ברזולוציה של $0.001\U{g}$.
  \item מיקרומטר, ברזולוציה של $0.001\U{mm}$.
  \item אריומטר, ברזולוציה של $0.001\U{g/cm^3}$.
  \item סרגל, ברזולוציה של $1\U{mm}$.
  \item מצלמת סמארטפון (\textenglish{Galaxy A12}) המצלמת 30 פריימים בשנייה.
  \item חצובה למצלמה.
  \item בריסטול לבן.
  \item חיישן טמפרטורה המחובר למחשב עם תוכנת \textenglish{Capstone}, ברזולוציה של $0.001\U{^\circ C}$.
  \item מחשב עם תוכנת העקיבה \textenglish{Tracker}.
  \item מגנט ומחט פרומגנטית: המחט ממוגנטת על ידי המגנט ומשמשת לאחיזת הכדורים, שקטנים מכדי לאחוז בהם באצבעות.
\end{enumerate}

% Schematic, Claude's drawing (no usable photo in the briefing). Side view: camera (8) on its tripod
% (9) faces the tube (1); bristol (10) behind the tube. LTR block: under RTL the x axis is mirrored.
\newcommand{\eqlabel}[3]{\node[draw, circle, fill=white, inner sep=1pt, font=\scriptsize] at (#1,#2) {#3};}
\begin{figure}[H]
  \centering
  \begin{otherlanguage}{english}
  \begin{tikzpicture}[x=1cm, y=1cm]
    % bristol behind the tube
    \fill[gray!8] (4.9,-0.3) rectangle (5.2,6.3); \draw (4.9,-0.3) rectangle (5.2,6.3);
    % tube with glycerin
    \fill[yellow!25] (3.2,0) rectangle (3.9,5.4);
    \draw[thick] (3.2,0) -- (3.2,5.8); \draw[thick] (3.9,0) -- (3.9,5.8); \draw[thick] (3.2,0) -- (3.9,0);
    \fill[gray!60] (3.15,5.8) rectangle (3.95,6.0); % valve / cap
    % height markers
    \fill[orange!80] (3.12,4.3) rectangle (3.98,4.45);
    \fill[orange!80] (3.12,1.0) rectangle (3.98,1.15);
    % ball and its direction of motion
    \fill[black!70] (3.55,3.2) circle (0.1);
    \draw[->, thick] (3.55,3.0) -- (3.55,2.4);
    \node[right, font=\scriptsize] at (3.6,2.7) {$v$};
    % camera on tripod
    \fill[black!80] (0.6,2.55) rectangle (1.0,3.35);
    \draw (0.8,2.55) -- (0.8,1.6); \draw (0.8,1.6) -- (0.4,0); \draw (0.8,1.6) -- (0.8,0); \draw (0.8,1.6) -- (1.2,0);
    \draw[dashed, gray] (1.0,3.25) -- (3.2,5.4); \draw[dashed, gray] (1.0,2.65) -- (3.2,0);
    % labels (numbers = equipment list)
    \draw[gray] (3.9,5.1) -- (4.3,5.1); \eqlabel{4.45}{5.1}{1}
    \draw[gray] (3.98,4.37) -- (4.3,4.37); \draw[gray] (3.98,1.07) -- (4.3,4.2); \eqlabel{4.45}{4.37}{2}
    \draw[gray] (3.65,3.2) -- (4.3,3.2); \eqlabel{4.45}{3.2}{3}
    \eqlabel{0.25}{3.0}{8} \eqlabel{1.25}{0.9}{9} \eqlabel{5.5}{3.0}{10}
  \end{tikzpicture}
  \end{otherlanguage}
  \caption{תרשים סכמטי של מערכת הניסוי, מבט מהצד: המצלמה על החצובה מול המבחנה המלאה בגליצרין, ומאחורי המבחנה הבריסטול. הקווים המקווקווים מסמנים את שדה הראייה של המצלמה. המספור תואם לרשימת הציוד.}
  \label{fig:system}
\end{figure}

% Photos of the setup, taken by the user in the lab (IMG_8522-8524, 2026-11-22; cropped, IMG_8522 rotated
% upright). Labels = equipment-list numbers, placed by Claude (fractions of the photo, origin bottom-left).
% Not in any photo: hydrometer (6), ruler (7), temperature sensor (11).
\newcommand{\labeledphoto}[3]{%
  \begin{tikzpicture}
    \node[anchor=south west, inner sep=0] (img) at (0,0) {\includegraphics[#1]{#2}};
    \begin{scope}[x={(img.south east)}, y={(img.north west)}] #3 \end{scope}
  \end{tikzpicture}}
\begin{figure}[H]
  \centering
  \begin{otherlanguage}{english}
  \labeledphoto{height=8.5cm}{figures/photo-tube.jpg}{%
    \eqlabel{0.45}{0.62}{1} \eqlabel{0.57}{0.72}{2} \eqlabel{0.57}{0.52}{2}
    \eqlabel{0.15}{0.70}{10} \eqlabel{0.85}{0.62}{12} \eqlabel{0.62}{0.10}{3}}%
  \hspace{0.4cm}%
  \labeledphoto{height=8.5cm}{figures/photo-camera.jpg}{%
    \eqlabel{0.30}{0.70}{8} \eqlabel{0.44}{0.17}{9} \eqlabel{0.88}{0.80}{1} \eqlabel{0.70}{0.86}{10}}
  \end{otherlanguage}
  \caption{תמונות של מערכת הניסוי במעבדה. משמאל: המבחנה עם סמני הגובה, הבריסטול מאחוריה, מגש הכדורים והמחשב. מימין: המצלמה על החצובה מול המבחנה. המספור תואם לרשימת הציוד.}
  \label{fig:photos-setup}
\end{figure}

\begin{figure}[H]
  \centering
  \begin{otherlanguage}{english}
  \labeledphoto{width=0.6\textwidth}{figures/photo-part-a.jpg}{%
    \eqlabel{0.17}{0.25}{4} \eqlabel{0.28}{0.88}{5} \eqlabel{0.83}{0.30}{13}}
  \end{otherlanguage}
  \caption{המכשירים לחלק א׳: המאזניים, המיקרומטר, והמגנט עם המחט הפרומגנטית המשמשת לאחיזת הכדורים. המספור תואם לרשימת הציוד.}
  \label{fig:photos-part-a}
\end{figure}

\subsection*{שלבי הניסוי}
\subsubsection*{חלק א׳ -- מדידת צפיפות הכדורים והגליצרין}
\begin{enumerate}
  \item צפיפות הגליצרין נמדדה באריומטר, שהוכנס לגליצרין וממנו נקראה הצפיפות ישירות מהסקאלה. המדידה בוצעה פעם אחת.
  \item לפי תדריך הניסוי, מכל אחת משש קבוצות הכדורים נבחרו באקראי שלושה כדורים (עם החזרה) מתוך כ-6 עד 10 הכדורים שבקבוצה. לכל כדור נמדדו המסה במאזניים הדיגיטליים והקוטר במיקרומטר.
  \item טמפרטורת הגליצרין נמדדה בחיישן הטמפרטורה ובתוכנת \textenglish{Capstone}, במחשב אחד לכל הקבוצות במעבדה, ונקבעה כממוצע הטמפרטורות שנמדדו לאורך 4 שעות, לפי תדריך הניסוי.
\end{enumerate}

\subsubsection*{חלק ב׳ -- מדידת מהירות נפילת הכדורים}
\begin{enumerate}
  \item המצלמה הוצבה על החצובה מול המבחנה, והבריסטול הלבן הוצב מאחורי המבחנה. הצילום נעשה בתמונה ריבועית (יחס צירים 1:1), כנדרש לפעולה תקינה של תוכנת העקיבה לפי תדריך הניסוי.
  \item שני סמני גובה הוצבו על המבחנה, והמרחק ביניהם נמדד פעם אחת בסרגל: $21\U{cm}$. מרחק זה שימש לכיול המרחקים בתוכנת העקיבה.
  \item מכל אחת משש הקבוצות הופל לתוך הגליצרין כדור אחד, ונפילתו צולמה בסרטון נפרד. הכדור נלקח מאותה קבוצה שממנה נבחרו הכדורים שנמדדו בחלק א׳.
  \item כל סרטון נותח בתוכנת \textenglish{Tracker}. תחילה כוילו המרחקים בסרטון: המרחק בין סמני הגובה הוגדר בתוכנה כ-$21\U{cm}$. לאחר מכן סומן הכדור בפריים שבו החלה העקיבה, והתוכנה עקבה אחריו בפריימים הבאים. בפריימים שבהם התוכנה לא זיהתה את הכדור נכון, מיקומו תוקן ידנית. בסרטונים של קבוצות 3 עד 6 הכיול לא עודכן, והמרחק בין הסמנים נותר על ערך ברירת המחדל של התוכנה, $1\U{m}$. לכן המיקומים בסרטונים אלה תוקנו בעיבוד הנתונים, לפי הנחיית מדריך הניסוי.
\end{enumerate}

\section*{תכנון עיבוד התוצאות}
% DRAFT 2026-11-22, Claude's wording (blue) for the user's review. Sources: briefing B p.3 (part a), p.6-7
% (steps 3-8), p.8 (appendix: why tau is not used), p.9 (appendix a, theoretical eta); general rules 3-4.4.
% Formula budget: 18 (theory (1)-(6) excluded). Numbered here: (7)-(16) = 10 of 18.
% Red \TODO = a choice the user has to make (no heuristic chosen silently).
\added{מטרת העיבוד היא לבחון את הקשר בין מהירות הנפילה הסופית לריבוע רדיוס הכדור, נוסחה \eqref{eq:vf}, ולחלץ ממנו את צמיגות הגליצרין. תחילה ייבחן אם המהירות הסופית תלויה לינארית בריבוע הרדיוס, ואם האיבר החופשי מתאפס. אם הקשר מתקיים, שיפוע ההתאמה מייצג את הביטוי $\frac{2g(\rho_b-\rho_l)}{9\eta}$ שבנוסחה \eqref{eq:vf}, ולכן ניתן לחלץ ממנו את הצמיגות ולהשוות אותה לערך התיאורטי. המהירות הסופית של כל כדור תחולץ בשתי דרכים, לפי תדריך הניסוי: בהתאמה א׳ ישירות מנתוני המהירות, ובהתאמה ב׳ מהתאמת נוסחה \eqref{eq:vt} למהירות כתלות בזמן. בשתי הדרכים תותאם לאחר מכן המהירות הסופית לריבוע הרדיוס, ומהשיפוע תחולץ הצמיגות. כל ההתאמות יבוצעו במחברת ה-\textenglish{Colab} של הקורס, בשיטת \textenglish{ODR}, שמקבלת את $x$, $\Delta x$, $y$, $\Delta y$. כל הגדלים יחושבו ביחידות \textenglish{c.g.s}.}

\subsection*{חלק א׳ -- צפיפות הכדורים והגליצרין}
\added{מסת כל כדור $m$ וקוטרו $d$ נמדדו פעם אחת, ולכן שגיאת כל אחד מהם היא שגיאת רזולוציה, המחושבת לפי נוסחה (28) בחוברת עיבוד הנתונים תחת הנחת התפלגות אחידה בין שנתות המכשיר, כאשר הרזולוציות הן של המאזניים ושל המיקרומטר. רדיוס הכדור הוא $r=\frac{d}{2}$, ושגיאתו $\Delta r=\frac{\Delta d}{2}$. נפח כל כדור יחושב לפי:}
\begin{equation}\label{eq:V}
V = \frac{4}{3}\pi r^3
\end{equation}
\added{ושגיאתו העקיפה, לפי נוסחה (39) בחוברת:}
\begin{equation}\label{eq:dV}
\Delta V = 4\pi r^2 \Delta r
\end{equation}

\added{צפיפות הכדורים תחולץ מהתאמה של המסה לנפח, לפי הקשר $m=\rho_b V$ (תדריך הניסוי, עמ׳ 3). מכיוון שכל הכדורים עשויים מאותו חומר, תותאם פונקציה אחת לכל 18 הכדורים שנמדדו. ההתאמה היא לינארית:}
\begin{equation}\label{eq:linfit}
y = a_0 + a_1 x
\end{equation}
\added{כאשר $x=V$ ו-$y=m$. הפרמטר $a_1$ מייצג את צפיפות הכדורים, והאיבר החופשי $a_0$ צפוי להתאפס. צפיפות הכדורים ושגיאתה יחולצו לפי:}
\begin{equation}\label{eq:rhob}
\rho_b = a_1 \pm \Delta a_1
\end{equation}
\added{צפיפות הגליצרין $\rho_l$ נקראה פעם אחת מהאריומטר, ושגיאתה היא שגיאת רזולוציה לפי נוסחה (28) בחוברת. טמפרטורת הגליצרין היא הערך שנמסר ואושר על ידי מדריך הניסוי, $T = 23.116 \pm 0.031\U{^\circ C}$.}

\added{הכדור שהופל מכל קבוצה לא נמדד בעצמו, אלא נלקח מאותה קבוצה שממנה נמדדו שלושה כדורים. לכן רדיוסו יילקח כממוצע שלושת הרדיוסים שנמדדו בקבוצה, $\bar r$. לרדיוס זה שני מקורות שגיאה: שגיאה סטטיסטית, מסטיית התקן המדגמית של שלוש המדידות (נוסחאות (32) ו-(34) בחוברת), ושגיאת הרזולוציה של המיקרומטר. השגיאה הכוללת מתקבלת לפי נוסחה (35) בחוברת:}
\begin{equation}\label{eq:drbar}
\Delta \bar r = \sqrt{\Delta r_{stat}^2 + \left(\frac{d_{res}}{2\sqrt{12}}\right)^2}
\end{equation}
\added{כאשר $\Delta r_{stat}$ היא השגיאה הסטטיסטית ו-$d_{res}$ רזולוציית המיקרומטר.}

\subsection*{חלק ב׳ -- נתוני הנפילה}
\added{לכל כדור התקבלו מתוכנת העקיבה הזמן $t_i$ והגובה $y_i$ בכל פריים $i$. בכדורים 3 עד 6 יוכפלו ערכי הגובה ב-21, לפי הנחיית מדריך הניסוי, כדי להמיר אותם לסנטימטרים (ראו שלבי הניסוי). מכל הזמנים יופחת הזמן של הפריים הראשון בעקיבה, כך שהזמן יימדד מתחילת העקיבה ולא מתחילת הסרטון (תדריך הניסוי, עמ׳ 7).}
\TODO{החלטה 1: שגיאת הזמן $\Delta t$ ושגיאת הגובה $\Delta y$ (התדריך, שלב 3, מבקש להעריך אותן). למשל: $\Delta t$ זניחה או שגיאת רזולוציה של מרווח פריים ($\frac{1}{30}\U{s}$), ו-$\Delta y$ כשגיאת רזולוציה של פיקסל אחד או כהערכת מיקום הכדור ב-\textenglish{Tracker}. מה בחרתם?}

\added{מהירות הנפילה בכל פריים תחושב כנגזרת נומרית של הגובה לפי הזמן, בגודלה (הכדור נע כלפי מטה):}
\begin{equation}\label{eq:vi}
v_i = \left|\frac{y_{i+1}-y_{i-1}}{t_{i+1}-t_{i-1}}\right|
\end{equation}
\added{ושגיאתה העקיפה, לפי נוסחה (43) בחוברת:}
\begin{equation}\label{eq:dvi}
\Delta v_i = \frac{\sqrt{\Delta y_{i+1}^2+\Delta y_{i-1}^2 + v_i^2\left(\Delta t_{i+1}^2+\Delta t_{i-1}^2\right)}}{t_{i+1}-t_{i-1}}
\end{equation}
\TODO{החלטה 2: נגזרת מרכזית (הנוסחה למעלה, בין הפריים שלפני לפריים שאחרי) או נגזרת קדמית (בין פריים לפריים שאחריו)? מרכזית מדויקת יותר אך מאבדת את הפריים הראשון והאחרון.}

\subsection*{התאמה א׳ -- המהירות הסופית מנתוני המהירות}
\added{לכל כדור תחולץ המהירות הסופית $v_f$ מנתוני המהירות בתחום שבו הכדור כבר הגיע למהירות קבועה.}
\TODO{החלטה 3: איך לחלץ את $v_f$ ומה שגיאתה (התדריך, שלב 4, שואל זאת)? אפשרות (א): ממוצע $v_i$ בתחום הקבוע, ושגיאה סטטיסטית (נוסחאות (32), (34) בחוברת). אפשרות (ב): התאמה לינארית (נוסחה \eqref{eq:linfit}) של $y$ כתלות ב-$t$ בתחום הקבוע, כך ש-$v_f=|a_1|$. בשתיהן צריך כלל דטרמיניסטי לבחירת התחום הקבוע. למשל: מהפריים שבו עבר זמן של כמה $\tau$ צפויים מתחילת הנפילה, לפי נוסחה \eqref{eq:vt}.}

\subsection*{התאמה ב׳ -- התאמה אקספוננציאלית למהירות}
\added{לכל כדור תותאם נוסחה \eqref{eq:vt} לנתוני המהירות כתלות בזמן, כאשר $x=t$ ו-$y=v$. הפרמטרים הם $v_f$, $v_0$ ו-$\tau$. $v_f$ צפויה להיות המהירות שאליה מתייצבים הנתונים בסוף הנפילה, $v_0$ המהירות בפריים הראשון, ו-$\tau$ זמן הדעיכה. אלה גם יהיו הניחושים ההתחלתיים להתאמה: $v_0$ כמהירות הראשונה שנמדדה, $v_f$ כמהירות האחרונה, ו-$\tau$ לפי נוסחה \eqref{eq:vt} עם הצמיגות התיאורטית. הצמיגות לא תחולץ מ-$\tau$: קצב הדגימה של המצלמה, כ-$30\U{ms}$, ארוך מזמן הדעיכה הצפוי, ונגזרת הגובה מגדילה את הרעש דווקא בפריימים הראשונים שקובעים את $\tau$ (תדריך הניסוי, נספח עמ׳ 8). לכן מהתאמה זו יילקח רק $v_f$, עם שגיאתו מההתאמה.}

\subsection*{חילוץ הצמיגות}
\added{בכל אחת משתי ההתאמות, המהירויות הסופיות של ששת הכדורים יותאמו לריבוע הרדיוס בהתאמה לינארית, נוסחה \eqref{eq:linfit}, כאשר $x=\bar r^2$ ו-$y=v_f$. השגיאה ב-$x$ היא שגיאה עקיפה לפי נוסחה (39) בחוברת:}
\begin{equation}\label{eq:dr2}
\Delta(\bar r^2) = 2\bar r\,\Delta \bar r
\end{equation}
\added{לפי נוסחה \eqref{eq:vf}, השיפוע $a_1$ מייצג את $\frac{2g(\rho_b-\rho_l)}{9\eta}$, והאיבר החופשי $a_0$ צפוי להתאפס. הצמיגות תחולץ לפי:}
\begin{equation}\label{eq:eta}
\eta = \frac{2g(\rho_b-\rho_l)}{9a_1}
\end{equation}
\added{ושגיאתה העקיפה, לפי נוסחה (43) בחוברת:}
\begin{equation}\label{eq:deta}
\Delta\eta = \eta\sqrt{\left(\frac{\Delta a_1}{a_1}\right)^2 + \frac{\Delta\rho_b^2+\Delta\rho_l^2}{(\rho_b-\rho_l)^2} + \left(\frac{\Delta g}{g}\right)^2}
\end{equation}
\added{כאשר $g$ תאוצת הכובד ו-$\Delta g$ שגיאתה.}
\TODO{החלטה 4: ערך $g$ ושגיאתו, ומאיזה מקור? אף קובץ של צוות הצמיגות לא נותן ערך. בדוח הנפילה החופשית נלקח $g = 981 \pm 10\U{cm/s^2}$ מקובץ התיאוריה של אותו ניסוי.}

\subsection*{בחינת ההתאמות}
\added{מכל התאמה יתקבלו הפרמטרים ושגיאותיהם, ולכל פרמטר תחושב השגיאה היחסית, באחוזים, לפי נוסחה (25) בחוברת. איכות כל התאמה תיבחן באמצעות המדדים הבאים:}
\begin{itemize}
  \item \added{חי בריבוע מצומצם, $\chi^2_{red}$ (נוסחה (50) בחוברת), המבטא את מידת ההתאמה בין נקודות המדידה לגרף ההתאמה. בהתאמה תקינה צפוי ערך קרוב ל-1.}
  \item \added{\textenglish{p-value}: ההסתברות לקבל ערך $\chi^2_{red}$ גדול מזה שהתקבל או שווה לו, בהנחה שהמודל מתאר את המערכת. ערכים תקינים הם בתחום $0.05 \le p \le 0.95$.}
  \item \added{גרף השארים, שבו ייבדק שפיזור הנקודות סביב קו האפס אקראי וללא מגמה.}
\end{itemize}
\added{הקשר הלינארי בין $v_f$ ל-$\bar r^2$ (מטרה 2) ייבחן באמצעות מדדים אלה בהתאמות של חילוץ הצמיגות. הצמיגות מכל אחת משתי ההתאמות תושווה לערך התיאורטי באמצעות מדד $N_\sigma$ (נוסחה (51) בחוברת). שתי התוצאות נחשבות תואמות אם $N_\sigma \le 3$.}
\TODO{החלטה 5: איך לקבל את הערך התיאורטי של $\eta$ ב-$T=23.116\U{^\circ C}$ מגרף נספח א׳ בתדריך (ביחידות סנטיפואז, $1\U{P}=100\U{cP}$)? אפשרות (א): קריאה ישירה מהגרף, עם שגיאת קריאה. אפשרות (ב): התאמת נוסחה \eqref{eq:etaT} מהרקע התיאורטי לנקודות המסומנות בגרף ($\times$ ב-6, 15, 20, 25, 30 מעלות) וחישוב ב-$T$. אפשרות (ג): אינטרפולציה לינארית בין הנקודות ב-20 וב-25 מעלות.}

\section*{תוצאות}
\TODO{תוצאות}

\section*{דיון}
\subsection*{עמידה במטרות}
\TODO{עמידה במטרות}
\subsection*{מקורות שגיאה}
\TODO{מקורות שגיאה}
\subsection*{תופעות שנצפו}
\TODO{תופעות שנצפו}
\subsection*{מגבלות והצעות לשיפור}
\TODO{מגבלות והצעות לשיפור}

\section*{שאלת מדריך}
\TODO{שאלת מדריך}

\section*{נספחים}
\TODO{נספחים}

\end{document}
